\section*{الفصل \ref{ch:b1:plate-kinematics} --- حركيّة الصفائح على الكرة}

\begin{solution}{exo:b1:plate-kinematics:1}
(أ) $\qty{0.210}{\degree\per\Myr} = 0.210\times\pi/180 = \qty{3.665e-3}{rad}$
لكل مليون سنة، أي \qty{3.665e-9}{rad.yr^{-1}}؛ ومع $\qty{1}{yr} = \qty{3.156e7}{s}$، يكون \qty{1.161e-16}{rad.s^{-1}}. (ب) $360/0.210 =
\qty{1714}{\Myr}$، أي أكثر من ثلث عمر الأرض. (ج) $\omega R =
\num{3.665e-9}\times\qty{6.371e9}{mm} = \qty{23.35}{mm.yr^{-1}}$.
\end{solution}

\begin{solution}{exo:b1:plate-kinematics:2}
$v = \omega R\sin\Delta = 23.35\sin\Delta$: \qty{11.7}{mm.yr^{-1}} عند \ang{30}، و\qty{20.2}{mm.yr^{-1}} عند \ang{60}، و\qty{23.35}{mm.yr^{-1}} عند \ang{90}. والانتشار أسرع ما يكون حيث يكون الحدّ أبعد عن القطب، حتى \ang{90}: بين أوراسيا وأمريكا الشمالية، عند الطرف الجنوبي من الحدّ، قرب الأزور، على بعد \ang{78} من القطب.
\end{solution}

\begin{solution}{exo:b1:plate-kinematics:3}
استمرّت الحقبة الفرعية $1.072 - 0.988 = \qty{0.084}{\Myr}$، والمعدّل النصفي بوحدة \unit{mm.yr^{-1}} هو أيضًا معدّل بوحدة \unit{km\per\Myr}. (أ) $50\times0.084 =
\qty{4.2}{km}$ عند \qty{50}{mm.yr^{-1}}؛ و$10\times0.084 = \qty{0.84}{km}$ عند
\qty{10}{mm.yr^{-1}}. (ب) الخط يعبر المحور: $2\times10\times0.781 = \qty{15.6}{km}$.
\end{solution}

\begin{solution}{exo:b1:plate-kinematics:4}
المعدّلان النصفيان هما $66/2.581 = \qty{25.6}{mm.yr^{-1}}$ غربًا و$64/2.581 = \qty{24.8}{mm.yr^{-1}}$ شرقًا؛ والمعدّل الكامل هو
$130/2.581 = \qty{50.4}{mm.yr^{-1}}$. الانتشار متماثل تقريبًا: فالجانب الغربي نما أسرع بنحو \qty{3}{\percent}.
\end{solution}

\begin{solution}{exo:b1:plate-kinematics:5}
$\cos\Delta = \sin\ang{61.8}\sin\ang{85} + \cos\ang{61.8}\cos\ang{85}
\cos(\ang{0} - \ang{139.6}) = 0.8813\times0.9962 + 0.4726\times0.08716\times
(-0.7615) = 0.87795 - 0.03137 = 0.84658$، فيكون $\Delta = \ang{32.16}$ و$v = 23.35\sin\ang{32.16} = \qty{12.4}{mm.yr^{-1}}$. وهذا أعلى قليلًا من \qtyrange{11}{12}{mm.yr^{-1}} المقاسة، لكن \omterm{def:b1:plate-kinematics:euler-pole}{متجهة أويلر} هي أفضل ملاءمة واحدة ل 453 معدّلًا، تبتعد عنها المعدّلات بمقدار \qty{0.7}{mm.yr^{-1}} في المتوسط، والمعدّلات الفردية قرب حيد غاكيل البطيء من بين الأقل دقة.
\end{solution}

\begin{solution}{exo:b1:plate-kinematics:6}
(أ) $\cos\Delta = \sin\ang{61.8}\sin\ang{39.44} + \cos\ang{61.8}\cos\ang{39.44}
\cos(\ang{169.43}) = 0.20111$: $\Delta = \ang{78.40}$. (ب) $v = 23.35
\sin\ang{78.40} = \qty{22.87}{mm.yr^{-1}}$، أي \qty{23}{mm.yr^{-1}} المقاسة
مباشرةً شمال المفصل. (ج) بصيغة الطريقة، البسط هو $\sin\ang{169.43}\cos\ang{61.8} = 0.08668$، والمقام
$\cos\ang{39.44}\sin\ang{61.8} - \sin\ang{39.44}\cos\ang{61.8}\cos\ang{169.43}
= 0.97573$: يُرى القطب عند $\beta = \ang{5.08}$، أي تقريبًا شمالًا، وتتحرك أوراسيا نحو $\ang{95.1}$. (د) الصدوع التحويلية موازية للحركة النسبية: تتجه نحو \ang{095}، أي تقريبًا شرق–غرب.
\end{solution}

\begin{solution}{exo:b1:plate-kinematics:7}
(أ) $\vec\omega = \omega(\cos\varphi_p\cos\lambda_p, \cos\varphi_p\sin\lambda_p,
\sin\varphi_p)$: $\vec\omega_{\mathrm{EU\,NA}} = (-0.0756, 0.0643, 0.1851)$ و
$\vec\omega_{\mathrm{NB\,NA}} = (0.0333, 0.0282, 0.2289)$، بوحدة
\unit{\degree\per\Myr}. (ب) $\vec\omega_{\mathrm{NB\,EU}} =
\vec\omega_{\mathrm{NB\,NA}} - \vec\omega_{\mathrm{EU\,NA}} = (0.1089, -0.0361,
0.0438)$، وطوله
\[
  \sqrt{0.011859 + 0.001303 + 0.001918} = \qty{0.1228}{\degree\per\Myr}.
\]
يقع قطبه عند \omterm{def:b1:spherical-astronomy:geographic}{العرض الجغرافي} $\arcsin(0.0438/0.1228) = \ang{20.9}$~ش \omterm{def:b1:spherical-astronomy:geographic}{والطول الجغرافي} $\arctan(-0.0361/0.1089) = \ang{-18.3}$، أي \ang{18.3}~غ (المركبة $x$ موجبة، فلا يُضاف \ang{180}). (ج) لا يقدّم الحدّ بين النوبية وأوراسيا إلا قياسات مباشرة قليلة: فبين الأزور وجبل طارق لا توجد إلا أسمات صدع غلوريا. ويُقاس الانتشاران عبر الأطلسي بمئات المعدّلات، ويرث فرقهما دقتهما.
\end{solution}

\begin{solution}{exo:b1:plate-kinematics:8}
(أ) حول $\vec e_z$ بزاوية \ang{90}: $\cos\ang{90} = 0$، و$\vec e_z\cdot\vec e_x = 0$
و$\vec e_z\times\vec e_x = \vec e_y$، إذن ينتقل $\vec e_x$ إلى $\vec e_y$. ثم حول $\vec e_x$ بزاوية \ang{90}: $\vec e_x\cdot\vec e_y = 0$ و
$\vec e_x\times\vec e_y = \vec e_z$، إذن ينتقل $\vec e_y$ إلى $\vec e_z$: فتنتهي النقطة عند القطب الشمالي. (ب) حول $\vec e_x$ أولًا: $\vec e_x$ على المحور، و$\vec e_x\times\vec e_x = \vec 0$ و$(1 - \cos\ang{90})(\vec e_x\cdot\vec e_x)
\vec e_x = \vec e_x$، إذن لا تتحرك النقطة؛ ثم حول $\vec e_z$ تنتقل إلى
$\vec e_y$، عند \ang{0}~ش و\ang{90}~ق. (ج) النتيجتان تختلفان: \omterm{def:b1:plate-kinematics:finite-rotation}{الدورانات المنتهية} لا تتبادل. وينبغي تركيب إعادات البناء المتعاقبة بضرب مصفوفات الدوران بالترتيب الصحيح، لا بجمع الزوايا أو الأقطاب أبدًا.
\end{solution}

\begin{solution}{exo:b1:plate-kinematics:9}
(أ) في فضاء السرع نضع $a$ عند الأصل، و$b$ على بعد \qty{30}{mm.yr^{-1}} نحو الشرق، و$c$ على بعد \qty{40}{mm.yr^{-1}} نحو الجنوب. (ب) $\vec v_{CB} =
\vec v_{CA} - \vec v_{BA} = (0, -40) - (30, 0) = (-30, -40)$ (مركبتا الشرق والشمال): \qty{50}{mm.yr^{-1}} نحو \omterm{def:b1:spherical-astronomy:horizontal}{السمت}
$\arctan\bigl((-30)/(-40)\bigr) + \ang{180} = \ang{216.9}$، أي جنوب غرب. (ج) الحيد عمودي على هذا الاتجاه، فيتّجه نحو
$\ang{216.9} - \ang{90} = \ang{126.9}$، أي شمال غرب–جنوب شرق.
\end{solution}

\begin{solution}{exo:b1:plate-kinematics:10}
(أ) تُسهم المصفوفة المحايدة بالمقدار $3\cos\alpha$، وتُسهم المصفوفة $\vec n\,\vec n^{\mathsf T}$ بأثرها $n_x^2 + n_y^2 + n_z^2 = 1$، مضروبًا في $1 - \cos\alpha$، أما المصفوفة المتناظرة المقلوبة فلا تُسهم بشيء: $\operatorname{tr}R =
3\cos\alpha + 1 - \cos\alpha = 1 + 2\cos\alpha$. (ب) هنا $\operatorname{tr}R = 0$،
إذن $\cos\alpha = -1/2$ و$\alpha = \ang{120}$. يحقّق المحور $R\vec n =
\vec n$: $R(x, y, z) = (z, x, y)$، إذن $x = y = z$ و$\vec n =
(1, 1, 1)/\sqrt3$، عند \omterm{def:b1:spherical-astronomy:geographic}{العرض الجغرافي} $\arcsin(1/\sqrt3) = \ang{35.26}$~ش والطول \ang{45}~ق. (ج) $R\vec e_x = \vec e_y$، $R\vec e_y = \vec e_z$،
$R\vec e_z = \vec e_x$: تتبدّل المحاور الثلاثة، كما يفعل ثلث دورة حول القطر. وصيغة رودريغ بالزاوية $\alpha = +\ang{120}$ حول $\vec n$ تعطي فعلًا
$\vec e_x\mapsto -\tfrac12\vec e_x + \tfrac12(1, 1, 1) +
\tfrac{\sqrt3}{2}\cdot\tfrac{1}{\sqrt3}(0, 1, -1) = (0, 1, 0)$، فالدوران عكس عقارب الساعة منظورًا من فوق \ang{35.26}~ش و\ang{45}~ق.
\end{solution}

\begin{solution}{exo:b1:plate-kinematics:11}
(أ) على صدع تحويلي تكون السرعة النسبية $\vec\omega\times\vec r$ على امتداد الصدع، أي موازية للمتجهة $\vec t$؛ وهي عمودية على $\vec\omega$، إذن $\vec t\cdot\vec\omega = 0$. والمستوي العمودي على $\vec t$ المار بالمركز يحوي $\vec r$ (لأن $\vec t$ أفقية) و$\vec\omega$: فهو مستوي \omterm{def:b1:spherical-astronomy:spherical-triangle}{الدائرة العظمى} المارّة بالنقطة والعمودية على الصدع، والقطب يقع على تلك الدائرة. ويتقاطع مستويان من هذا النوع على امتداد $\vec t_1\times\vec t_2$. (ب) $\vec t_1 = (0.5436, 0.8367, -0.0673)$ عند المفصل، و$\vec t_2 = (0.4942, 0.8639, -0.0968)$ في أيسلندا، إذن $\vec t_1\times\vec t_2 = (-0.02285, 0.01936, 0.05612)$، وطوله $0.06362$. عرضه الجغرافي $\arcsin(0.05612/0.06362) = \ang{61.9}$~ش، وطوله الجغرافي $\ang{180} - \arctan(0.01936/0.02285) = \ang{139.7}$~ق (المركبة $x$ سالبة): هو قطب الفصل. (ج) بصدعين تحويليين وحدهما يُحدَّد القطب تحديدًا رديئًا: فالدائرتان العظميان تتقاطعان بزاوية صغيرة، لأن الاتجاهين يختلفان بمقدار \ang{8} فقط على امتداد \ang{26} من الحدّ. وتغيير أيّ من الاتجاهين بدرجة واحدة يحرّك القطب بنحو \ang{15}. والتحديدات الحقيقية تجمع صدوعًا تحويلية كثيرة مع معدّلات الانتشار، إذ يحدّد تغيّرها على طول الحدّ البعد الزاوي إلى القطب.
\end{solution}

\begin{solution}{exo:b1:plate-kinematics:12}
(أ) $(94 - 4)/3 = \qty{30}{mm.yr^{-1}}$. (ب) الفروق بين موضعين يبعد أحدهما عن الآخر سنة هي
$24 - 4 = 20$، و$25 - 5 = 20$، و$26 - 6 = 20$، و$65 - 15 = 50$، و$74 - 24 = 50$،
$75 - 25 = 50$، و$76 - 26 = 50$، و$85 - 65 = 20$ و$94 - 74 = 20$ (بالمليمتر). مرتبةً: خمس قيم هي 20 وأربع هي 50؛ والوسيط \qty{20}{mm.yr^{-1}}. (ج) موضعان يبعد أحدهما عن الآخر سنة يقعان في الطور نفسه من التذبذب الموسمي، فيُلغى ذلك في فرقهما. والقفزة تؤثر فقط في الأزواج التي تعبرها، وهنا أربعة من التسعة؛ وما دامت أقل من نصف الأزواج تعبر قفزة، يتجاهلها الوسيط، بينما يأخذ التقدير بين الطرفين في (أ) القفزة كلها حركةً فيخرج أعلى بنسبة \qty{50}{\percent}. ويفشل الأسلوب حين تكون القفزات من التكرار بحيث تحتوي معظم أزواج السنة الواحدة على واحدة منها، أو حين تتحرك المحطة بلا انتظام، كما بعد زلزال كبير.
\end{solution}

\begin{solution}{pb:b1:plate-kinematics:1}
\textbf{1.} $\omega = 0.210\times(\pi/180)/\num{e6} = \qty{3.665e-9}{rad.yr^{-1}}$؛
$\omega R = \num{3.665e-9}\times\qty{6.371e9}{mm} = \qty{23.35}{mm.yr^{-1}}$.

\textbf{2.} $\cos\Delta = \sin\ang{61.8}\sin\ang{64.5} + \cos\ang{61.8}
\cos\ang{64.5}\cos(\ang{-157.6}) = 0.79546 - 0.18809 = 0.60737$، إذن
$\Delta = \ang{52.60}$.

\textbf{3.} $v = 23.35\sin\ang{52.60} = 23.35\times0.7944 =
\qty{18.55}{mm.yr^{-1}}$.

\textbf{4.} البسط هو $\sin\ang{157.6}\cos\ang{61.8} = 0.18011$، والمقام
\[
  \cos\ang{64.5}\sin\ang{61.8} - \sin\ang{64.5}\cos\ang{61.8}\cos\ang{157.6} =
  0.37940 + 0.39436 = 0.77376;
\]
كلاهما موجب، إذن $\beta = \arctan 0.2328 =
\ang{13.1}$. وتتحرك أوراسيا نحو $\ang{13.1} + \ang{90} = \ang{103.1}$.

\textbf{5.} شرقًا: $18.55\sin\ang{103.1} = \qty{18.07}{mm.yr^{-1}}$؛ شمالًا:
$18.55\cos\ang{103.1} = \qty{-4.20}{mm.yr^{-1}}$، أي \qty{4.20}{mm.yr^{-1}} نحو الجنوب.

\textbf{6.} $\cos\Delta = 0.8851$، $\Delta = \ang{27.74}$، $v = 23.35
\sin\ang{27.74} = \qty{10.9}{mm.yr^{-1}}$. تدور الصفيحتان حول قطب في شرق سيبيريا: فحيد غاكيل قريب منه، وتتناقص السرعة مع $\sin\Delta$ رغم أن الصفيحتين هما نفسهما.

\textbf{7.} غربًا: $94/9.779 = \qty{9.61}{mm.yr^{-1}}$؛ شرقًا: $97/9.779 =
\qty{9.92}{mm.yr^{-1}}$.

\textbf{8.} $(94 + 97)/9.779 = 191/9.779 = \qty{19.53}{mm.yr^{-1}}$.

\textbf{9.} من C2An: $54.1/2.581 = \qty{20.96}{mm.yr^{-1}}$؛ ومن قاعدة C1n: $19.8/0.781 = \qty{25.35}{mm.yr^{-1}}$. فهما لا يتفقان: فكلما كان الحدّ أحدث، كان المعدّل أعلى. قد يكون التغيّر الحقيقي في المعدّل ممكنًا، لكن زيادة ثابتة في المسافات، متساوية في كل الأعمار، تُحدث هذا بالضبط.

\textbf{10.} نسبة المسافة غربًا إلى المسافة شرقًا تتراوح بين 0.96 و0.98: فالجانب الشرقي نما بنسبة بضعة في المئة أكثر. ويعتمد المعدّل الكامل على مجموع المسافتين وحده، فلا يغيّره عدم التماثل.

\textbf{11.} $(191 - 5.0)/9.779 = \qty{19.02}{mm.yr^{-1}}$؛
$(54.1 - 5.0)/2.581 = \qty{19.02}{mm.yr^{-1}}$؛ $(19.8 - 5.0)/0.781 =
\qty{18.95}{mm.yr^{-1}}$. والثلاثة تتفق الآن في حدود \qty{0.1}{mm.yr^{-1}}. فالإزاحة نفسها \qty{5}{km} هي ربع \qty{19.8}{km} للزوج الأحدث، لكنها أقل من \qty{3}{\percent} من \qty{191}{km} للزوج الأقدم؛ والحدود القديمة تعطي معدّلات أكثر موثوقية.

\textbf{12.} المعدّل النصفي المصحّح هو $186/9.779/2 = \qty{9.51}{mm.yr^{-1}}$،
و$9.51\times0.084 = \qty{0.80}{km}$، أي ثلث \qty{2.5}{km} التي يُزاح بها كل حدّ. ويُتوقَّع أن يُطمس خط بهذا الضيق. ولأنه مُتعرَّف عليه على حيد ريكيانيس، بينما هو غائب عن حيد غاكيل الفقير بالصهارة، فذلك يوحي بأن الإمداد الوافر بالصهارة يسجّل المجال بأمانة حتى في الخطوط الضيقة.

\textbf{13.} المعدّلات المصحّحة، من \qty{18.95}{mm.yr^{-1}} إلى \qty{19.02}{mm.yr^{-1}}، تتفق مع \qty{19.0}{mm.yr^{-1}} لمتجهة الصفيحة؛ ومن دون التصحيح تزيد عليها بنسبة 3 إلى 33 في المئة.

\textbf{14.} $\cos\Delta = \sin\ang{64.139}\sin\ang{59.018} +
\cos\ang{64.139}\cos\ang{59.018}\cos\ang{27.554} = 0.97054$، فيكون $\Delta =
\ang{13.94}$؛ والوتر هو $2R\sin(\Delta/2) = 12742\times\sin\ang{6.97} =
\qty{1546}{km}$.

\textbf{15.} $1000/17.5 = 57$ سنة.

\textbf{16.} يختلف الاتجاهان بمقدار $\ang{123.1} - \ang{114.8} =
\ang{8.3}$: المركبة الأفقية على \omterm{def:b1:spherical-astronomy:spherical-triangle}{الدائرة العظمى} هي
$18.9\cos\ang{8.3} = \qty{18.70}{mm.yr^{-1}}$، وعلى الوتر
$18.70\cos\ang{6.97} = \qty{18.56}{mm.yr^{-1}}$، أي نحو \qty{18.6}{mm.yr^{-1}}.

\textbf{17.} يطول الوتر بمقدار \qty{17.5}{mm.yr^{-1}}، أي أقل بنسبة \qty{6}{\percent} من \qty{18.6}{mm.yr^{-1}} المتنبأ بها. وتقع المحطتان بعيدًا عن حدّ الصفيحة، فالفرق ليس محليًا. فإما أن الصفيحتين تتحركان اليوم أبطأ قليلًا مما تحركتا به في المتوسط خلال الثلاثة ملايين سنة الأخيرة، وإما أن أحد التقديرين يحمل خطأً منهجيًا صغيرًا؛ والقياسات وحدها لا تستطيع أن تحسم ذلك. \omterm{def:b1:plate-kinematics:euler-pole}{ومتجهة أويلر} المحسوبة من تحديد المواقع بالأقمار الصناعية وحدها تتنبأ بمقدار \qty{17.1}{mm.yr^{-1}} لهذا الوتر، وهي قريبة من القيمة المقاسة.

\textbf{18.} شرقًا: $21.12 - (-1.07) = \qty{22.19}{mm.yr^{-1}}$؛ شمالًا:
$-6.97 - (-0.97) = \qty{-6.00}{mm.yr^{-1}}$. السرعة $\sqrt{22.19^2 + 6.00^2} =
\qty{23.0}{mm.yr^{-1}}$، نحو $\ang{180} - \arctan(22.19/6.00) =
\ang{105.1}$.

\textbf{19.} الاتجاه قريب من \ang{103} لمتجهة الصفيحة، لكن السرعة تزيد على \qty{18.6}{mm.yr^{-1}} بربعها. ويقع المحطتان داخل النطاق الذي يتشوّه فيه الحدّ؛ وهوفن، قرب أكبر غطاء جليدي، تتحرك مع ارتداد \omterm{def:b1:layered-earth:layers}{القشرة} مع ترقّق الجليد: إلى الأعلى بمقدار \qty{12.9}{mm.yr^{-1}}، وجانبيًا، بعيدًا عن الجليد. والمحطات داخل منطقة الحدّ تقيس حركة الصفيحة مضافًا إليها التشوّه المحلي.

\textbf{20.} $0.210\times10 = \ang{2.10}$.

\textbf{21.} تسير النقطة على \omterm{def:b1:spherical-astronomy:spherical-triangle}{دائرة صغرى} نصف قطرها $R\sin\Delta$:
$6371\times0.7944\times(2.10\times\pi/180) = \qty{185.5}{km}$. والعرض المصحّح \omterm{def:b1:layered-earth:layers}{للقشرة} المتكوّنة منذ الحافة الحديثة للحقبة C5n، \qty{186}{km} في \qty{9.8}{Myr}، قريب منه: فالحركة كانت ثابتة طيلة عشرة ملايين سنة.
\textbf{22.} المركبتان (شرقًا، شمالًا): أوراسيا
$(22.9\sin\ang{95.1}, 22.9\cos\ang{95.1}) = (22.81, -2.04)$، والنوبية
$(19.2\sin\ang{103.8}, 19.2\cos\ang{103.8}) = (18.65, -4.58)$. أوراسيا بالنسبة
إلى النوبية: $(4.16, 2.54)$، أي \qty{4.9}{mm.yr^{-1}} نحو
$\arctan(4.16/2.54) = \ang{58.6}$، شرق شمال شرق.

\textbf{23.} مفصل من ثلاثة حيود مستقرّ أيًّا كانت السرع. والسرعة في السؤال 22 تجعل أوراسيا والنوبية يبتعدان، من شرق شمال شرق إلى غرب جنوب غرب: فيجب أن ينفتح حدّهما عند المفصل، ببطء؛ وهذا هو شقّ تيرسيرا، واتجاهه شرق جنوب شرق، فيجعل الانفتاح مائلًا.

\textbf{24.} الجزء الأول: \qty{18.6}{mm.yr^{-1}} في وسط أيسلندا (\qty{19.0}{mm.yr^{-1}} عند قطاع الجزء الثاني). الجزء الثاني: من 18.95 إلى \qty{19.02}{mm.yr^{-1}} بعد التصحيح. الجزء الثالث: يطول الوتر بين ريكيافيك وستافانغر بمقدار \qty{17.5}{mm.yr^{-1}}، بينما تتنبأ متجهة الصفيحة بمقدار \qty{18.6}{mm.yr^{-1}}؛ والزوج داخل أيسلندا يعطي \qty{23.0}{mm.yr^{-1}}، وهو ليس حركة الصفيحة. وتتفق متجهة الصفيحة مع الخطوط المغناطيسية؛ وتختلف الأقمار الصناعية بنسبة \qty{6}{\percent}. وتتوسّط الخطوط على ما يصل إلى عشرة ملايين سنة، وتتوسّط متجهة الصفيحة على \qty{3.16}{Myr}؛ أما الأقمار الصناعية فتقيس الحركة في آخر عشرين إلى ثلاثين سنة.

\textbf{25.} تعطي متجهة الصفيحة \qty{18.55}{mm.yr^{-1}} عبر أيسلندا، والخطوط المصحّحة \qty{19.0}{mm.yr^{-1}} بجانبها: بأقرب مليمتر في السنة، فإن \omterm{def:b1:plate-kinematics:spreading-rate}{معدّل الانتشار الكامل} عبر أيسلندا هو
\textbf{\qty[text-series-to-math=true]{19}{mm.yr^{-1}}}.
\end{solution}
